\documentclass[../ap_2021.tex]{subfiles}

\graphicspath{{./image/}}

\begin{document}

\setcounter{chapter}{3}
\chapter{電磁気学:双極子放射}
\section{}
\begin{equation}\begin{aligned}[b]
    \varphi(\vb*{r})&=\frac{q}{4\pi\varepsilon_0\abs{\vb*{r}-\frac{\vb*{l}}{2}}}-\frac{q}{4\pi\varepsilon_0\abs{\vb*{r}+\frac{\vb*{l}}{2}}}\\
    &\simeq \frac{q}{4\pi\varepsilon_0r(1-\frac{l}{r}\cos\theta)^{1/2}}-\frac{q}{4\pi\varepsilon_0r(1+\frac{l}{r}\cos\theta)^{1/2}}\\
    &\simeq \frac{q}{4\pi\varepsilon_0r}\qty(1+\frac{l}{2r}\cos\theta)-\frac{q}{4\pi\varepsilon_0r}\qty(1-\frac{l}{2r}\cos\theta)\\
    &= \frac{\alpha_0E_0\cos\theta}{4\pi\varepsilon_0r^2}
\end{aligned}\end{equation}
より\(C=0,D=1\)

\section{}
\begin{equation}\begin{aligned}[b]
    \varphi(\vb*{r},t)
    &=\frac{q(t-\frac{r_+}{c})}{4\pi\varepsilon_0r_+}-\frac{q(t-\frac{r_-}{c})}{4\pi\varepsilon_0r_-}\\
    &\simeq \frac{q(t-\frac{r_+}{c})}{4\pi\varepsilon_0r}\qty(1+\frac{l}{2r}\cos\theta)-\frac{q(t-\frac{r_+}{c})}{4\pi\varepsilon_0r}\qty(1-\frac{l}{2r}\cos\theta)\\
    &\simeq \frac{1}{4\pi\varepsilon_0 r}\qty[q\qty(r-\frac{r}{c}+\frac{l\cos\theta}{2cr})-q\qty(r-\frac{r}{c}-\frac{l\cos\theta}{2cr})]
        +\frac{\cos\theta}{4\pi\varepsilon_0r^2}\qty(\frac{p(t-\frac{r_+}{c})+p(t-\frac{r_-}{c})}{2})\\
    &\simeq\frac{1}{4\pi\varepsilon_0r}\frac{\dot{q}(t-\frac{r}{c})l\cos\theta}{c}+\frac{\cos\theta}{4\pi\varepsilon_0r^2}p\qty(t-\frac{r}{c})\\
    &=\frac{\alpha_0E_0\cos\theta}{4\pi\varepsilon_0}\qty{-\frac{\omega_0}{c}\sin(\omega_0\qty(t-\frac{r}{c}))\cdot\frac{1}{r}
    +\cos(\omega_0\qty(t-\frac{r}{c}))\cdot\frac{1}{r^2}}
\end{aligned}\end{equation}
より
\begin{equation}\begin{aligned}[b]
    F &= -\frac{\omega_0}{c}\sin(\omega_0\qty(t-\frac{r}{c})), & G = \cos(\omega_0\qty(t-\frac{r}{c}))
\end{aligned}\end{equation}

\section{}
\begin{equation}\begin{aligned}[b]
    \vb*{A}(\vb*{r},t) = -\frac{\alpha_0E_0}{4\pi\varepsilon_0 cr} \frac{\omega_0}{c}\sin(\omega_0\qty(t-\frac{r}{c}))\qty{\cos\theta\vb*{e}_r-\sin\theta\vb*{e}_\theta}
\end{aligned}\end{equation}

\section{}
勾配を取る操作は\(r\)方向以外必ず\(r^{-2}\)になり、
\(r\)方向も\(\sin(\omega_0(t-\frac{r}{c}))\)以外を微分しても\(r^{-2}\)になる。なので
\begin{equation}\begin{aligned}[b]
    \grad{\varphi}(\vb*{r},t)=\frac{\alpha_0E_0\cos\theta}{4\pi\varepsilon_0r}\qty(\frac{\omega_0}{c})^2\cos(\omega_0\qty(t-\frac{r}{c}))\vb*{e}_r
\end{aligned}\end{equation}

\section{}
\(\vb*{A}(\vb*{r},t)\)の時間微分は
\begin{equation}\begin{aligned}[b]
    \pdv{\vb*{A}(\vb*{r},t)}{t}=-\frac{\alpha_0E_0}{4\pi\varepsilon_0 r} \qty(\frac{\omega_0}{c})^2\cos(\omega_0\qty(t-\frac{r}{c}))\qty{\cos\theta\vb*{e}_r-\sin\theta\vb*{e}_\theta}
\end{aligned}\end{equation}

これより電場は
\begin{equation}\begin{aligned}[b]
    \vb*{E}(\vb*{r},t)
    &=-\frac{\alpha_0E_0\cos\theta}{4\pi\varepsilon_0r}\qty(\frac{\omega_0}{c})^2\cos(\omega_0\qty(t-\frac{r}{c}))\vb*{e}_r
        +\frac{\alpha_0E_0}{4\pi\varepsilon_0 r} \qty(\frac{\omega_0}{c})^2\cos(\omega_0\qty(t-\frac{r}{c}))\qty{\cos\theta\vb*{e}_r-\sin\theta\vb*{e}_\theta}\\
    &= -\frac{\alpha_0E_0\sin\theta\omega_0^2}{4\pi\varepsilon_0c^2}\frac{\cos(\omega_0\qty(t-\frac{r}{c}))}{r}\vb*{e}_\theta
\end{aligned}\end{equation}

磁場について、勾配の計算と同じような点に気を付けると
\begin{equation}\begin{aligned}[b]
    \vb*{B}(\vb*{r},t)&=\curl \vb*{A}(\vb*{r},t) = \frac{1}{r}\pdv{r}(rA_\theta)\vb*{e}_\phi\\
    &= -\frac{\mu_0\alpha_0E_0\sin\theta \omega_0^2}{4\pi c}\frac{\cos(\omega_0\qty(t-\frac{r}{c}))}{r}\vb*{e}_\phi
\end{aligned}\end{equation}

\section{}
\begin{equation}\begin{aligned}[b]
    \vb*{S}(\vb*{r},t)=\frac{1}{\mu_0}\vb*{E}(\vb*{r},t)\times \vb*{B}(\vb*{r},t)
    = \frac{\alpha_0^2E_0^2\sin^2\theta \omega^4}{16\pi^2\varepsilon_0c^3}\frac{\cos^2\qty(\omega_0\qty(t-\frac{r}{c}))}{r^2}\vb*{e}_r
\end{aligned}\end{equation}
図は略

\section{}
\begin{equation}\begin{aligned}[b]
    \vb*{E}(\vb*{r},t)&=-\frac{\alpha(t-\frac{r}{c})E_0\sin\theta\omega_0^2}{4\pi\varepsilon_0c^2}\frac{\cos(\omega_0\qty(t-\frac{r}{c}))}{r}\vb*{e}_\theta\\
    &=-\frac{\alpha_0E_0\sin\theta\omega_0^2}{4\pi\varepsilon_0c^2}\frac{\cos(\omega_0\qty(t-\frac{r}{c}))}{r}\vb*{e}_\theta
        -\frac{\alpha_1E_0\sin\theta\omega_0^2}{4\pi\varepsilon_0c^2}\frac{\cos(\omega_v\qty(t-\frac{r}{c}))\cos(\omega_0\qty(t-\frac{r}{c}))}{r}\vb*{e}_\theta\\
    &=-\frac{\alpha_0E_0\sin\theta\omega_0^2}{4\pi\varepsilon_0c^2}\frac{\cos(\omega_0\qty(t-\frac{r}{c}))}{r}\vb*{e}_\theta\\
        &\qquad-\frac{\alpha_1E_0\sin\theta\omega_0^2}{8\pi\varepsilon_0c^2}\frac{\cos((\omega_0-\omega_v)\qty(t-\frac{r}{c}))}{r}\vb*{e}_\theta
        -\frac{\alpha_1E_0\sin\theta\omega_0^2}{8\pi\varepsilon_0c^2}\frac{\cos((\omega_0+\omega_v)\qty(t-\frac{r}{c}))}{r}\vb*{e}_\theta\\
\end{aligned}\end{equation}


\section*{感想}
やってるときには苦しかったけど、
答案を見ると実は分量少なくて驚いた。
双極子放射は慣れると結構早く計算できるものなのかもしれないね。

最後の問題はラマン散乱を見てるなぁと分かったので、
周波数は\(\omega_0\pm\omega_1\)になると予想がついた。

\(\vb*{A}\)の時間微分のところに出てきそうな気はしたが、
もし、これも考慮して微分すると\(\omega_v\ll\omega_0\)の係数が外に出てきてしまうので、
結局考えなくても問題ないと分かる。
それで計算を進めると確かにラマンっぽい電場になって安心。


\end{document}
